% ============================================================
% WS1 reframing text — drop-in replacements for sn-article.tex
% Numbers marked [DMRG: ...] are placeholders to be filled from
% dmrg_refs/*.npz once the validation runs complete.
% ============================================================

% ---------------------------------------------------------------
% (1) REPLACE the second half of Introduction paragraph 1
% (currently: "...while classical numerical simulations rapidly
% become computationally intractable as system size increases.")
% A referee will object that in 1D this is false (DMRG). New text
% concedes this up front and states the actual gap.
% ---------------------------------------------------------------

Although recent advances in programmable quantum simulators and quantum
computers have enabled the preparation of increasingly complex quantum
states \cite{bourgund2025formation, king2025beyond}, obtaining and
analyzing large-scale quantum measurements remains prohibitively expensive
in physical experiments. On the classical side, the applicability of
numerical methods is regime-dependent: in one dimension, tensor-network
methods such as the density-matrix renormalization group (DMRG) provide
quasi-exact ground-state properties at modest cost, whereas in higher
dimensions, for long-range or frustrated interactions, and for states with
extensive entanglement, their cost grows rapidly. Most importantly, all
such methods require exact knowledge of the Hamiltonian, whereas
present-day quantum hardware produces measurement records from devices
whose effective Hamiltonians are known only approximately. Developing
scalable approaches that infer large-scale quantum behavior directly from
limited, experimentally accessible measurement data therefore remains a
fundamental open challenge.

% ---------------------------------------------------------------
% (2) ADD to the final Introduction paragraph (after "...including
% the transverse-field Ising model, the frustrated J1-J2 model, and
% the axial next-nearest-neighbor Ising (ANNNI) model."):
% ---------------------------------------------------------------

We deliberately benchmark in one dimension, where quasi-exact
tensor-network references remain available at every system size we
generate: this allows the cross-scale generation and extrapolation
pipeline to be validated rigorously, rather than only checked for internal
consistency. The one-dimensional study is thus a controlled proof of
principle for a method whose intended regime is precisely where such
classical references cease to exist---higher-dimensional and frustrated
systems, and measurement data from quantum hardware.

% ---------------------------------------------------------------
% (3) NEW Results paragraph — insert at the end of
% "Thermodynamic-limit inference from recursively generated
% measurements" (fill numbers from dmrg_refs).
% ---------------------------------------------------------------

Because exact diagonalization is limited to $N\lesssim 24$, we validate
the recursively generated measurements at larger sizes against DMRG
calculations performed at the same system sizes and couplings
($N=32$--$96$, periodic boundary conditions, bond dimension $\chi=256$,
energy convergence $10^{-9}$; the DMRG implementation reproduces exact
diagonalization at $N=12$ to machine precision). Away from the critical
region, observables estimated from the recursively generated measurements
agree with DMRG to a mean absolute deviation of $0.017$ (all distances;
$0.025$ for $r=1$) across all sizes up to $N=96$ --- four times the largest
exactly diagonalizable system. Deviations concentrate near the critical
point, where correlation lengths are largest: for $g\in[0.4,0.6]$ the mean
$r{=}1$ deviation grows to $0.086$ (maximum $0.167$ at $N=96$, $g=0.5$),
increasing systematically with recursion depth, and the generated
correlations decay too rapidly with distance on the paramagnetic side.
The recursive procedure is therefore quantitatively reliable in both
phases but accumulates bias in the critical region, and we restrict
thermodynamic-limit claims accordingly. The $1/N$-extrapolated thermodynamic-limit values agree with independent
infinite-system (iDMRG) estimates to better than $0.01$ deep in the
ferromagnetic phase ($g\le0.2$) and at $g=0.8$; in the critical region the
extrapolations inherit the recursion bias and undershoot the
infinite-system values by up to $0.13$, and deep in the paramagnet a small
residual offset ($\approx0.05$) survives extrapolation. This analysis
converts the consistency-across-rounds
check of the recursive procedure (Supplementary Fig.~4) into a validation
against quasi-exact references at every size reached by the recursion ---
including a transparent map of where the procedure degrades.

% ---------------------------------------------------------------
% (4) REPLACE the Discussion sentence "...Diffushadow provides a
% data-driven pathway that complements existing analytical and
% numerical approaches..." with an explicit relation-to-tensor-
% networks paragraph:
% ---------------------------------------------------------------

It is useful to state precisely how this framework relates to established
tensor-network methods. For the one-dimensional ground states studied
here, DMRG is superior in both accuracy and cost, and we use it as the
validation standard at system sizes beyond exact diagonalization. The
advantages of learned measurement generation lie on orthogonal axes.
First, its input is measurement data rather than a Hamiltonian: it applies
directly to records produced by quantum hardware, for which the effective
Hamiltonian---and hence any variational simulation---is only approximately
known. Second, a single trained model is amortized across the full
parameter range and across system sizes, whereas each tensor-network
calculation addresses one parameter point at one size. Third, its
per-sample generation cost is insensitive to the entanglement structure of
the underlying state, while tensor-network costs grow exponentially with
lattice width in two dimensions; our two-dimensional results demonstrate
the mechanism precisely in this direction. Finally, generated measurements
are unlimited in number at fixed cost, making sample-hungry estimators
practical, although generative amplification reduces statistical variance
only and cannot remove model bias. We therefore regard learned measurement
generation not as a competitor to tensor networks in their domain of
strength, but as a complementary inference route whose natural targets are
experimental measurement data and regimes beyond one dimension.

% ---------------------------------------------------------------
% (5) ADD near the end of the Discussion (Rydberg bridge + outlook):
% ---------------------------------------------------------------

The ANNNI benchmark also places Diffushadow in the universality class of
commensurate--incommensurate transitions realized by Rydberg-atom arrays
\cite{ebadi2021quantum, chepiga2019floating}, whose native output---fixed-basis
occupation snapshots across a scanned detuning---is exactly the data format
consumed by our framework. Applying measurement-based generative inference
to such experimental records, for instance to locate incommensurate and
floating-phase boundaries at system sizes beyond those directly
characterized, is a natural next step. % (bib entries to add: ebadi2021quantum
% = Ebadi et al., Nature 595, 227 (2021); chepiga2019floating = Chepiga &
% Mila, Phys. Rev. Lett. 122, 017205 (2019).)
